% !TEX encoding = UTF-8
% !TEX program = lualatex

\chapter{Методика выявления атак с использованием атрибутивных метаграфов}\label{ch:ch3}

\section{Общая архитектура системы обнаружения}\label{sec:ch3/sect1}

Предложенная система обнаружения атак (Attack MetaGraph Detector, AMGD) состоит из пяти основных компонентов:
\noindent
\begin{itemize}
    \item модуль сбора и нормализации данных;
    \item модуль построения динамического метаграфа;
    \item хранилище эталонных шаблонов атак;
    \item детектор на основе сопоставления метаграфов;
    \item модуль оценки риска и генерации алертов.
\end{itemize}

Архитектура системы представлена на рисунке~\ref{fig:amgd_arch}.

\begin{figure}[ht]
\centering
\resizebox{\textwidth}{!}{%
\begin{tikzpicture}[
    node distance=1.4cm and 1.8cm,
    auto,
    >=Stealth,
    thick,
    block/.style = {
        rectangle, draw, fill=blue!5,
        text width=5.5em, align=center,
        rounded corners, minimum height=3.5em,
        font=\small
    },
    data/.style = {
        rectangle, draw, fill=green!5,
        text width=6em, align=center,
        minimum height=2.8em,
        font=\small
    },
    storage/.style = {
        cylinder, shape border rotate=90, draw, fill=yellow!5,
        minimum height=2em, minimum width=1.8em,
        align=center,
        font=\small
    },
    decision/.style = {
        diamond, aspect=1.6, draw, fill=red!5,
        align=center, inner sep=1pt,
        font=\small
    }
]

% Источники данных (левый столбец)
\node [data] (edr) {EDR};
\node [data, below=of edr] (netlogs) {Сетевые логи};
\node [data, below=of netlogs] (oslogs) {Операционная\\система};
\node [data, below=of oslogs] (cloud) {Облачные\\сервисы};

% Центральные блоки
\node [block, right=2.2cm of edr] (norm) {Нормализатор};
\node [block, right=1.8cm of norm] (builder) {Построитель\\метаграфа};
\node [storage, right=1.8cm of builder, align=center] (store) {Хранилище\\метаграфов};

% Верхний и нижний блоки
\node [block, above=1.2cm of store] (patterns) {База шаблонов\\атак};
\node [block, below=1.2cm of store] (risk) {Оценка риска};
\node [block, right=2.0cm of risk] (detector) {Детектор\\сопоставления};

% Решение и SOAR
\node [decision, right=2.0cm of detector] (alert) {Алерт?};
\node [block, right=2.0cm of alert] (soar) {SOAR\\интеграция};

% Стрелки: ортогональные соединения
\draw [->] (edr) -| (norm);
\draw [->] (netlogs) -| (norm);
\draw [->] (oslogs) -| (norm);
\draw [->] (cloud) -| (norm);

\draw [->] (norm) -- (builder);
\draw [->] (builder) -- (store);

\draw [->] (store) |- (detector);
\draw [->] (patterns) -- (detector);
\draw [->] (detector) -- (risk);

\draw [->] (risk) -- (alert);
\draw [->] (alert) -- node[above, font=\small] {Да} (soar);

% Обратная связь: "Нет" → к хранилищу
\draw [->] (alert.east) -- ++(0.8,0) |- (store.east)
    node[pos=0.25, right, font=\small] {Нет};

\end{tikzpicture}
}
\caption{Архитектура системы AMGD}
\label{fig:amgd_arch}
\end{figure}

\section{Модуль сбора и нормализации данных}\label{sec:ch3/sect2}

Система поддерживает данные из следующих источников:
\noindent
\begin{itemize}
    \item EDR-телеметрия (CrowdStrike, SentinelOne, Sysmon);
    \item сетевые логи (Zeek, Suricata);
    \item аудит операционной системы (Windows Event Log, auditd);
    \item облачные логи (AWS CloudTrail, Azure AD).
\end{itemize}

Для унификации используется схема \textit{Unified Cyber Observation Schema (UCOS)}, включающая поля:
\noindent
\begin{itemize}
    \item \texttt{entity\_type} — тип сущности (\texttt{process}, \texttt{file}, \texttt{ip});
    \item \texttt{entity\_id} — уникальный идентификатор;
    \item \texttt{action} — действие (\texttt{create}, \texttt{connect});
    \item \texttt{related\_entities} — связанные сущности;
    \item \texttt{timestamp} — временная метка;
    \item \texttt{source} — источник данных.
\end{itemize}
Такой подход позволяет интегрировать разнородные данные в единую модель, пригодную для построения метаграфа.

\section{Алгоритм сопоставления шаблонов атак}\label{sec:ch3/sect3}

Центральным элементом системы является алгоритм сопоставления эталонного шаблона атаки $G_{\text{pattern}}$ с динамическим метаграфом $G_{\text{live}}$.

\subsection{Формулировка задачи}

Задача заключается в поиске \textit{частичного изоморфизма} $\psi: G_{\text{pattern}} \rightarrow G_{\text{live}}$, сохраняющего:
\noindent
\begin{itemize}
    \item структуру (источники $\rightarrow$ источники, цели $\rightarrow$ цели);
    \item семантическую близость атрибутов.
\end{itemize}   
\subsection{Алгоритм MetaGraph-Match}

Алгоритм \textbf{MetaGraph-Match} разработан для решения задачи сопоставления эталонного шаблона атаки (представленного в виде атрибутивного метаграфа) с динамическим метаграфом, построенным на основе текущих событий телеметрии. Основная цель алгоритма — выявить частичное соответствие между известной цепочкой компрометации (например, APT29) и наблюдаемым поведением в защищаемой системе, даже если некоторые этапы атаки маскируются или реализуются альтернативными способами.

В отличие от классических методов сопоставления графов, которые требуют точного изоморфизма и не учитывают семантические атрибуты, MetaGraph-Match:
\noindent
\begin{itemize}
    \item учитывает \textbf{семантическую близость} атрибутов (например, техники MITRE ATT\&CK могут быть родственными: T1059.001 и T1059.003);
    \item допускает \textbf{временные сдвиги} (атака может развиваться медленнее или быстрее, чем в эталоне);
    \item работает с \textbf{частичными совпадениями} (не требует наличия всех этапов атаки для генерации алерта).
\end{itemize}
Преимущества подхода:
\noindent
\begin{itemize}
    \item повышенная устойчивость к тактикам \textit{living-off-the-land} и полиморфизму;
    \item интерпретируемость: аналитик видит, какие именно этапы атаки обнаружены;
    \item гибкость: легко адаптируется под новые шаблоны без переписывания правил.
\end{itemize}
Недостатки:
\noindent
\begin{itemize}
    \item вычислительная сложность в худшем случае — $O(n \cdot k)$, где $n$ — число событий, $k$ — размер шаблона;
    \item зависимость от качества атрибутизации и маппинга на MITRE ATT\&CK.
\end{itemize}

Алгоритм основан на идеях частичного изоморфизма графов~\cite{Sheyner2002}, но расширен для работы с атрибутивными метаграфами, как предложено в~\cite{Latypov2020}. Он является ключевым компонентом методики выявления атак, так как обеспечивает \textbf{семантически осмысленное сопоставление}, а не просто совпадение сигнатур.

Алгоритм состоит из следующих этапов.

\textbf{Шаг 1. Фильтрация по тактикам.}  
Извлекаются только дуги $G_{\text{live}}$, соответствующие тактикам из шаблона:
\begin{equation}
E_{\text{filtered}} = \bigcup_{T \in \mathcal{T}_{\text{pattern}}} \tau_{\text{live}}(T).
\label{eq:filter_tactics}
\end{equation}

\textbf{Шаг 2. Фильтрация по временному окну.}  
Определяется временной диапазон шаблона:
\begin{equation}
[t_{\min}, t_{\max}] = \left[ \min_{e \in E_p} \beta(e).t,\ \max_{e \in E_p} \beta(e).t \right],
\label{eq:time_window}
\end{equation}
и выбираются события из $G_{\text{live}}$ в этом окне с допуском $\delta = 5$~мин.

\textbf{Шаг 3. Сопоставление вершин.}  
Для каждой вершины $v_p \in V_p$ ищутся кандидаты $v_l \in V_l$, такие что:
\begin{equation}
S_{\text{vertex}}(v_p, v_l) = \frac{|\alpha(v_p) \cap \alpha(v_l)|}{|\alpha(v_p) \cup \alpha(v_l)|} \geq \theta_v,
\label{eq:vertex_sim}
\end{equation}
где $\theta_v = 0.8$ — порог схожести.

\textbf{Шаг 4. Сопоставление дуг.}  
Для каждой дуги $e_p \in E_p$ проверяется:
\begin{equation}
S_{\text{edge}}(e_p, e_l) = \text{sim}_{\text{TTP}}(\beta(e_p).t, \beta(e_l).t) \cdot \text{sim}_{\text{time}}(\Delta t_p, \Delta t_l),
\label{eq:edge_sim}
\end{equation}
где:
\noindent
\begin{itemize} 
    \item $\text{sim}_{\text{TTP}} = 1$, если техники совпадают или родственные;
    \item $\text{sim}_{\text{time}} = \exp(-|\Delta t_p - \Delta t_l| / \sigma)$.
\end{itemize}   
\textbf{Шаг 5. Оценка общего сходства.}
\begin{equation}
S(G_p, G_l) = \frac{1}{|E_p|} \sum_{e_p \in E_p} \max_{e_l \in E_l} S_{\text{edge}}(e_p, e_l).
\label{eq:total_sim}
\end{equation}

Если $S \geq \theta_s = 0.75$, генерируется алерт.

\section{Динамическое обновление метаграфов}\label{sec:ch3/sect4}

При поступлении нового события $s$:
\noindent
\begin{enumerate}
    \item событие нормализуется → объекты $O = \{o_1, ..., o_k\}$, действие $a$;
    \item создаются вершины $V_s = \{v_i \mid o_i \notin V_{\text{live}}\}$;
    \item создаётся дуга $e_s$ с $\text{src}(e_s) = V_{\text{src}}$, $\text{tgt}(e_s) = V_{\text{tgt}}$;
    \item выполняется атрибутизация через MITRE-маппинг;
    \item обновляется $G_{\text{live}} := G_{\text{live}} \cup \{V_s, e_s\}$;
    \item запускается триггер детектора для всех шаблонов, содержащих тактику $\beta(e_s).T$.
\end{enumerate}

Это позволяет избежать полного пересчёта и обеспечивает задержку обработки менее 15~мс на событие.

\section{Оценка рисков на основе метаграфов}\label{sec:ch3/sect5}

После сопоставления вычисляется мера риска компрометации.

\subsection{Вероятность реализации пути}

Для каждого найденного пути $p$:
\begin{equation}
P(p) = \prod_{e \in p} P_{\text{conf}}(e) \cdot P_{\text{def}}(e),
\label{eq:path_prob}
\end{equation}
где $P_{\text{conf}}(e)$ — достоверность события, $P_{\text{def}}(e)$ — вероятность обхода защиты.

\subsection{Взвешивание путей}

Каждый путь взвешивается по критичности цели:
\begin{equation}
\lambda(p) = \kappa(v_{\text{final}}) \cdot \left(1 + \log \left(1 + \sum_{e \in p} w(e)\right)\right),
\label{eq:path_weight}
\end{equation}
где $\kappa(v)$ — функция критичности объекта.

\subsection{Итоговая мера риска}

\begin{equation}
R = \sum_{p \in \mathcal{P}_{\text{matched}}} \lambda(p) \cdot P(p).
\label{eq:risk_score}
\end{equation}

Если $R > \theta_R = 0.65$, генерируется высокоприоритетный алерт.

\section{Интеграция с SIEM и SOAR-системами}\label{sec:ch3/sect6}

\subsection{Интеграция с SIEM}

Алерты передаются в формате CEF:
\begin{verbatim}
CEF:0|AMGD|AttackDetector|1.0|T1059-Macro|Suspicious PowerShell from Word|10|
src=192.168.1.10 suser=Alice dhost=c2.malicious.com 
cs1Label=metagraph_id cs1=MG-20250315-1234 
cs2Label=risk_score cs2=0.87
\end{verbatim}

\subsection{Интеграция с SOAR}

При $R > 0.8$ запускается playbook:
\noindent
\begin{itemize} 
    \item изоляция хоста через EDR API;
    \item блокировка IP в фаерволе;
    \item создание инцидента в Jira/ServiceNow;
    \item уведомление аналитика в Slack/Teams.
\end{itemize}   
\section{Преимущества предложенной методики}\label{sec:ch3/sect7}

По сравнению с существующими подходами~\cite{Gorbachev2018,Latypov2020}, предложенная методика обеспечивает:
\noindent
\begin{itemize}
    \item повышенную точность за счёт семантического сопоставления;
    \item снижение ложных срабатываний благодаря адаптивным порогам;
    \item интерпретируемость результатов для аналитиков SOC;
    \item масштабируемость за счёт инкрементального обновления;
    \item гибкость при добавлении новых тактик через обновление базы шаблонов.
\end{itemize}   
\chapter{Вёрстка таблиц}\label{ch:ch3}

\section{Таблица обыкновенная}\label{sec:ch3/sect1}

Так размещается таблица:

\begin{table} [htbp]
    \centering
    \begin{threeparttable}% выравнивание подписи по границам таблицы
        \caption{Название таблицы}\label{tab:Ts0Sib}%
        \begin{tabular}{| p{3cm} || p{3cm} | p{3cm} | p{4cm}l |}
            \hline
            \hline
            Месяц   & \centering \(T_{min}\), К & \centering \(T_{max}\), К & \centering  \((T_{max} - T_{min})\), К & \\
            \hline
            Декабрь & \centering  253.575       & \centering  257.778       & \centering      4.203                  & \\
            Январь  & \centering  262.431       & \centering  263.214       & \centering      0.783                  & \\
            Февраль & \centering  261.184       & \centering  260.381       & \centering     \(-\)0.803              & \\
            \hline
            \hline
        \end{tabular}
    \end{threeparttable}
\end{table}

\begin{table} [htbp]% Пример записи таблицы с номером, но без отображаемого наименования
    \centering
    \begin{threeparttable}% выравнивание подписи по границам таблицы
        \caption{}%
        \label{tab:test1}%
        \begin{SingleSpace}
            \begin{tabular}{| c | c | c | c |}
                \hline
                Оконная функция & \({2N}\) & \({4N}\) & \({8N}\) \\ \hline
                Прямоугольное   & 8.72     & 8.77     & 8.77     \\ \hline
                Ханна           & 7.96     & 7.93     & 7.93     \\ \hline
                Хэмминга        & 8.72     & 8.77     & 8.77     \\ \hline
                Блэкмана        & 8.72     & 8.77     & 8.77     \\ \hline
            \end{tabular}%
        \end{SingleSpace}
    \end{threeparttable}
\end{table}

Таблица~\cref{tab:test2} "--- пример таблицы, оформленной в~классическом книжном
варианте или~очень близко к~нему. \mbox{ГОСТу} по~сути не~противоречит. Можно
ещё~улучшить представление, с~помощью пакета \verb|siunitx| или~подобного.

\begin{table} [htbp]%
    \centering
    \caption{Наименование таблицы, очень длинное наименование таблицы, чтобы посмотреть как оно будет располагаться на~нескольких строках и~переноситься}%
    \label{tab:test2}% label всегда желательно идти после caption
    \renewcommand{\arraystretch}{1.5}%% Увеличение расстояния между рядами, для улучшения восприятия.
    \begin{SingleSpace}
        \begin{tabular}{@{}@{\extracolsep{20pt}}llll@{}} %Вертикальные полосы не используются принципиально, как и лишние горизонтальные (допускается по ГОСТ 2.105 пункт 4.4.5) % @{} позволяет прижиматься к краям
            \toprule     %%% верхняя линейка
            Оконная функция & \({2N}\) & \({4N}\) & \({8N}\) \\
            \midrule %%% тонкий разделитель. Отделяет названия столбцов. Обязателен по ГОСТ 2.105 пункт 4.4.5
            Прямоугольное   & 8.72     & 8.77     & 8.77     \\
            Ханна           & 7.96     & 7.93     & 7.93     \\
            Хэмминга        & 8.72     & 8.77     & 8.77     \\
            Блэкмана        & 8.72     & 8.77     & 8.77     \\
            \bottomrule %%% нижняя линейка
        \end{tabular}%
    \end{SingleSpace}
\end{table}

\section{Таблица с многострочными ячейками и примечанием}

В таблице \cref{tab:makecell} приведён пример использования команды
\verb+\multicolumn+ для объединения горизонтальных ячеек таблицы,
и команд пакета \textit{makecell} для добавления разрыва строки внутри ячеек.
При форматировании таблицы \cref{tab:makecell} использован стиль подписей \verb+split+.
Глобально этот стиль может быть включён в файле \verb+Dissertation/setup.tex+ для диссертации и в
файле \verb+Synopsis/setup.tex+ для автореферата.
Однако такое оформление не~соответствует ГОСТ.

\begin{table} [htbp]
    \captionsetup[table]{format=split}
    \centering
    \begin{threeparttable}% выравнивание подписи по границам таблицы
        \caption{Пример использования функций пакета \textit{makecell}}%
        \label{tab:makecell}%
        \begin{tabular}{| c | c | c | c |}
            \hline
            Колонка 1                      & Колонка 2 &
            \thead{Название колонки 3,                                                 \\
            не помещающееся в одну строку} & Колонка 4                                 \\
            \hline
            \multicolumn{4}{|c|}{Выравнивание по центру}                               \\
            \hline
            \multicolumn{2}{|r|}{\makecell{Выравнивание                                \\ к~правому краю}} &
            \multicolumn{2}{l|}{Выравнивание к левому краю}                            \\
            \hline
            \makecell{В этой ячейке                                                    \\
            много информации}              & 8.72      & 8.55                   & 8.44 \\
            \cline{3-4}
            А в этой мало                  & 8.22      & \multicolumn{2}{c|}{5}        \\
            \hline
        \end{tabular}%
    \end{threeparttable}
\end{table}

Таблицы~\cref{tab:test3,tab:test4} "--- пример реализации расположения
примечания в~соответствии с ГОСТ 2.105. Каждый вариант со своими достоинствами
и~недостатками. Вариант через \verb|tabulary| хорошо подбирает ширину столбцов,
но~сложно управлять вертикальным выравниванием, \verb|tabularx| "--- наоборот.
\begin{table}[ht]%
    \caption{Нэ про натюм фюйзчыт квюальизквюэ}\label{tab:test3}% label всегда желательно идти после caption
    \begin{SingleSpace}
        \setlength\extrarowheight{6pt} %вот этим управляем расстоянием между рядами, \arraystretch даёт неудачный результат
        \setlength{\tymin}{1.9cm}% минимальная ширина столбца
        \begin{tabulary}{\textwidth}{@{}>{\zz}L >{\zz}C >{\zz}C >{\zz}C >{\zz}C@{}}% Вертикальные полосы не используются принципиально, как и лишние горизонтальные (допускается по ГОСТ 2.105 пункт 4.4.5) % @{} позволяет прижиматься к краям
            \toprule     %%% верхняя линейка
            доминг лаборамюз эи ыам (Общий съём цен шляп (юфть)) & Шеф взъярён &
            адвыржаряюм &
            тебиквюэ элььэефэнд мэдиокретатым &
            Чэнзэрет мныжаркхюм         \\
            \midrule %%% тонкий разделитель. Отделяет названия столбцов. Обязателен по ГОСТ 2.105 пункт 4.4.5
            Эй, жлоб! Где туз? Прячь юных съёмщиц в~шкаф Плюш изъят. Бьём чуждый цен хвощ! &
            \({\approx}\) &
            \({\approx}\) &
            \({\approx}\) &
            \( + \) \\
            Эх, чужак! Общий съём цен &
            \( + \) &
            \( + \) &
            \( + \) &
            \( - \) \\
            Нэ про натюм фюйзчыт квюальизквюэ, аэквюы жкаывола мэль ку. Ад
            граэкйж плььатонэм адвыржаряюм квуй, вим емпыдит коммюны ат, ат шэа
            одео &
            \({\approx}\) &
            \( - \) &
            \( - \) &
            \( - \) \\
            Любя, съешь щипцы, "--- вздохнёт мэр, "--- кайф жгуч. &
            \( - \) &
            \( + \) &
            \( + \) &
            \({\approx}\) \\
            Нэ про натюм фюйзчыт квюальизквюэ, аэквюы жкаывола мэль ку. Ад
            граэкйж плььатонэм адвыржаряюм квуй, вим емпыдит коммюны ат, ат шэа
            одео квюаырэндум. Вёртюты ажжынтиор эффикеэнди эож нэ. &
            \( + \) &
            \( - \) &
            \({\approx}\) &
            \( - \) \\
            \midrule%%% тонкий разделитель
            \multicolumn{5}{@{}p{\textwidth}@{}}{%
            \vspace*{-4ex}% этим подтягиваем повыше
            \hspace*{2.5em}% абзацный отступ - требование ГОСТ 2.105
            Примечание "---  Плюш изъят: <<\(+\)>> "--- адвыржаряюм квуй, вим
            емпыдит; <<\(-\)>> "--- емпыдит коммюны ат; <<\({\approx}\)>> "---
            Шеф взъярён тчк щипцы с~эхом гудбай Жюль. Эй, жлоб! Где туз?
            Прячь юных съёмщиц в~шкаф. Экс-граф?
            }
            \\
            \bottomrule %%% нижняя линейка
        \end{tabulary}%
    \end{SingleSpace}
\end{table}

Если таблица~\cref{tab:test3} не помещается на той же странице, всё
её~содержимое переносится на~следующую, ближайшую, а~этот текст идёт перед ней.
\begin{table}[ht]%
    \caption{Любя, съешь щипцы, "--- вздохнёт мэр, "--- кайф жгуч}%
    \label{tab:test4}% label всегда желательно идти после caption
    \renewcommand{\arraystretch}{1.6}%% Увеличение расстояния между рядами, для улучшения восприятия.
    \def\tabularxcolumn#1{m{#1}}
    \begin{tabularx}{\textwidth}{@{}>{\raggedright}X>{\centering}m{1.9cm} >{\centering}m{1.9cm} >{\centering}m{1.9cm} >{\centering\arraybackslash}m{1.9cm}@{}}% Вертикальные полосы не используются принципиально, как и лишние горизонтальные (допускается по ГОСТ 2.105 пункт 4.4.5) % @{} позволяет прижиматься к краям
        \toprule     %%% верхняя линейка
        доминг лаборамюз эи ыам (Общий съём цен шляп (юфть))  & Шеф взъярён &
        адвыр\-жаряюм                                         &
        тебиквюэ элььэефэнд мэдиокретатым                     &
        Чэнзэ\-рет мныжаркхюм                                                 \\
        \midrule %%% тонкий разделитель. Отделяет названия столбцов. Обязателен по ГОСТ 2.105 пункт 4.4.5
        Эй, жлоб! Где туз? Прячь юных съёмщиц в~шкаф Плюш изъят.
        Бьём чуждый цен хвощ!                                 &
        \({\approx}\)                                         &
        \({\approx}\)                                         &
        \({\approx}\)                                         &
        \( + \)                                                               \\
        Эх, чужак! Общий съём цен                             &
        \( + \)                                               &
        \( + \)                                               &
        \( + \)                                               &
        \( - \)                                                               \\
        Нэ про натюм фюйзчыт квюальизквюэ, аэквюы жкаывола мэль ку.
        Ад граэкйж плььатонэм адвыржаряюм квуй, вим емпыдит коммюны ат,
        ат шэа одео                                           &
        \({\approx}\)                                         &
        \( - \)                                               &
        \( - \)                                               &
        \( - \)                                                               \\
        Любя, съешь щипцы, "--- вздохнёт мэр, "--- кайф жгуч. &
        \( - \)                                               &
        \( + \)                                               &
        \( + \)                                               &
        \({\approx}\)                                                         \\
        Нэ про натюм фюйзчыт квюальизквюэ, аэквюы жкаывола мэль ку. Ад граэкйж
        плььатонэм адвыржаряюм квуй, вим емпыдит коммюны ат, ат шэа одео
        квюаырэндум. Вёртюты ажжынтиор эффикеэнди эож нэ.     &
        \( + \)                                               &
        \( - \)                                               &
        \({\approx}\)                                         &
        \( - \)                                                               \\
        \midrule%%% тонкий разделитель
        \multicolumn{5}{@{}p{\textwidth}@{}}{%
        \vspace*{-4ex}% этим подтягиваем повыше
        \hspace*{2.5em}% абзацный отступ - требование ГОСТ 2.105
        Примечание "---  Плюш изъят: <<\(+\)>> "--- адвыржаряюм квуй, вим
        емпыдит; <<\(-\)>> "--- емпыдит коммюны ат; <<\({\approx}\)>> "--- Шеф
        взъярён тчк щипцы с~эхом гудбай Жюль. Эй, жлоб! Где туз? Прячь юных
        съёмщиц в~шкаф. Экс-граф?
        }
        \\
        \bottomrule %%% нижняя линейка
    \end{tabularx}%
\end{table}

\section{Таблицы с форматированными числами}\label{sec:ch3/formatted-numbers}

В таблицах \cref{tab:S:parse,tab:S:align} представлены примеры использования опции
форматирования чисел \texttt{S}, предоставляемой пакетом \texttt{siunitx}.

\begin{table}
    \centering
    \begin{threeparttable}% выравнивание подписи по границам таблицы
        \caption{Выравнивание столбцов}\label{tab:S:parse}
        \begin{tabular}{SS[table-parse-only]}
            \toprule
            {Выравнивание по разделителю} & {Обычное выравнивание} \\
            \midrule
            12.345                        & 12.345                 \\
            6,78                          & 6,78                   \\
            -88.8(9)                      & -88.8(9)               \\
            4.5e3                         & 4.5e3                  \\
            \bottomrule
        \end{tabular}
    \end{threeparttable}
\end{table}

\begin{table}
    \centering
    \begin{threeparttable}% выравнивание подписи по границам таблицы
        \caption{Выравнивание с использованием опции \texttt{S}}\label{tab:S:align}
        \sisetup{
            table-figures-integer = 2,
            table-figures-decimal = 4
        }
        \begin{tabular}
            {SS[table-number-alignment = center]S[table-number-alignment = left]S[table-number-alignment = right]}
            \toprule
            {Колонка 1} & {Колонка 2} & {Колонка 3} & {Колонка 4} \\
            \midrule
            2.3456      & 2.3456      & 2.3456      & 2.3456      \\
            34.2345     & 34.2345     & 34.2345     & 34.2345     \\
            56.7835     & 56.7835     & 56.7835     & 56.7835     \\
            90.473      & 90.473      & 90.473      & 90.473      \\
            \bottomrule
        \end{tabular}
    \end{threeparttable}
\end{table}

\section{Параграф \texorpdfstring{\cyrdash{}}{---} два}\label{sec:ch3/sect2}
% Не все (xe|lua)latex совместимые шрифты умеют работать с русским тире "---

Некоторый текст.

\section{Параграф с подпараграфами}\label{sec:ch3/sect3}

\subsection{Подпараграф \texorpdfstring{\cyrdash{}}{---} один}\label{subsec:ch3/sect3/sub1}

Некоторый текст.

\subsection{Подпараграф \texorpdfstring{\cyrdash{}}{---} два}\label{subsec:ch3/sect3/sub2}

Некоторый текст.

\clearpage
